% Source: https://artofproblemsolving.com/community/c3h2757952_algebra_problem_with_3_variables
% Stable identity: aops-topic-2757952-q01
We have $x+2y=kz$, $y+2z=kx$ and $z+2x=ky$. Adding gives $3(x+y+z)=k(x+y+z)$.

If $x+y+z\ne0$, then $k=3$, which occurs (take $x=y=z=1$).

If $x+y+z=0$, substitute $z=-x-y$ into the second equation: $y-2x-2y=kx$, so $y=-(k+2)x$. The first equation then gives $x+2y=k(-x-y)$, i.e. $(1+k)x+(2+k)y=0$, and so $x\bigl[(1+k)-(k+2)^2\bigr]=0$. Since $x\ne0$ (otherwise $y=z=0$), $k^2+3k+3=0$, which has no real roots.

So the only value is $\boxed{k=3}$. (Over the complex numbers, the two roots $\dfrac{-3\pm i\sqrt3}2$ would also occur, and the values would sum to $0$.)
